%------------------------------------------------------------------------------%
\section{Potências de números negativos}

Vimos acima que é possível definir $a^x$ para todo $a>0$ e $x\in\R$.

\emph{Por que não definimos para $a\leq 0$?} \\
\emph{Faz sentido falar em $a^x$ para $a\leq0$?}

A resposta é positiva se $x$ é um número inteiro.
Se $x$ não é inteiro, essa definição não pode ser feita sem gerar
inconsistências, ou seja, não é possível definir.
Por exemplo,
\[
    (-4)^{\sfrac{1}{2}}
\]
não está definido, pois a definição natural para esse número seria
\[
  \sqrt{-4}
\]
que não é um número real.

Aparentemente, poderíamos definir
\[
  (-8)^{\sfrac{1}{3}}
\]
como
\[
  \sqrt[3]{-8} = -2.
\]
Mas há um problema que surge da equivalência de frações.
Sabemos que
\[
  \frac{1}{3} = \frac{2}{6},
\]
logo teríamos que
\[
  (-8)^{\sfrac{1}{3}}
  = (-8)^{\sfrac{2}{6}}
\]
mas
\[
  (-8)^{\sfrac{2}{6}}
  = \sqrt[6]{(-8)^2}
  = \sqrt[6]{2^6}
  = 2.
\]
Ou seja, uma potência de base negativa e expoente racional não pode ser
definida. Consequentemente, a definição de $a^x$ para $x$ irracional também não
está definida para $a<0$. Portanto
\[
  (-2)^{\sfrac{1}{2}}
  \quad\text{e}\quad
  (-3)^{\sqrt{3}}
\]
não são números reais.
